\section{Thread-Level Parallelism: Multithreading, SMT, Multicore}\label{ch:tlp}

\subsection{Decision map}
\small
\begin{tabularx}{\textwidth}{L{4.4cm}L{4.2cm}X}
\toprule
\tblhead{\textbf{Scenario} & \textbf{Feature} & \textbf{Why}}\\
\midrule
Throughput server, many independent threads/requests & Multicore (CMP) with multithreading per core & Applications with TLP (DB, web) map directly; ILP alone is saturated. \\
Pipeline stalls are \emph{long and rare} (e.g.\ L2 misses), want no slowdown of a single thread & Coarse-grained multithreading & Switch only on costly stalls; pipeline refill $\ll$ stall time. \\
Stalls are \emph{short and frequent}; hardware can switch every cycle & Fine-grained multithreading & Hides short and long stalls; interleaves each cycle. \\
Out-of-order superscalar with idle issue slots; want cheap extra throughput & SMT & Reuses renaming and OoO hardware; fills unused slots from another thread. \\
Single important thread must run as fast as possible & One big (SMT-capable) superscalar core; preferred-thread priority & Multi-core does not speed one thread; SMT with priority protects it. \\
Limited die area / thermal budget but threaded software & Several small cores & Each core simpler and easier to design/manufacture; better energy per throughput. \\
Applications hard to parallelise & Do \emph{not} rely on cores/threads & TLP gives no gain; spend on ILP/caches. \\
\bottomrule
\end{tabularx}
\normalsize

\subsection{Why TLP at all}
\begin{itemize}
\item ILP gives implicit parallelism within a loop or basic block; it saturates at 3--6 issue (Section~\ref{ch:ilp}), and for most applications \textbf{most execution units sit idle} even on an 8-way superscalar (Tullsen et al.).
\item \textbf{TLP is explicit:} multiple threads, each with its own instructions, data, PC and register state. Goals: improve \emph{throughput} of machines running many programs and \emph{execution time} of multithreaded programs. TLP ``could be more cost-effective to exploit than ILP''.
\item Natural parallelism: online transaction systems (many independent queries/updates), data mining, databases, scientific codes; server-per-client threads; games (AI, graphics, physics as threads).
\item Industry move: IBM Power4 (2000) was the first commercial single-chip general-purpose multiprocessor (2 Power3 + shared L2); Sun, AMD, Intel followed.
\end{itemize}
\src{ca06 ``Performance beyond single thread ILP'', ``Thread Level Parallelism (TLP)''; ca08 ``Why multi-core?'', ``Thread-level parallelism''}

\subsection{Multithreading on one core}
Multiple threads \textbf{share the functional units} of one processor by overlapping. The processor \emph{duplicates the independent state of each thread}: a separate register file, a separate PC, and (for independent programs) a separate page table. Memory is shared via the virtual-memory mechanism. Hardware thread switch is \emph{much faster} than a full process switch (100s--1000s of clocks).

\begin{center}
\small
\begin{tabular}{L{2.4cm}L{4.2cm}L{4.5cm}L{4.4cm}}
\toprule
\tblhead{\textbf{} & \textbf{Fine-grained} & \textbf{Coarse-grained} & \textbf{SMT (simultaneous)}}\\
\midrule
When it switches & Every instruction/clock, round-robin, skipping stalled threads & Only on costly stalls (e.g.\ L2 miss) & No switching: instructions from several threads issue in the \emph{same cycle} \\
Hides & Short \emph{and} long stalls & Only long stalls (pipeline refill $\ll$ stall time) & Short and long stalls \emph{and} wasted issue slots \\
Single-thread cost & Slows a ready thread (delayed by others) & Does not slow a thread (others run only when it stalls) & Can slow a thread unless priority is given \\
Hardware & Switch every cycle; per-thread state & No very fast switch needed; pipeline must be flushed/frozen on switch & OoO superscalar + per-thread rename table, PC, (logically) per-thread ROB \\
Example & Sun Niagara & IBM AS/400 & Pentium 4 (HT), IBM Power5 \\
\bottomrule
\end{tabular}
\end{center}

\begin{dec}[Pick the multithreading style]
\begin{itemize}
\item \textbf{Coarse-grained} if the core is simple/in-order, stalls are dominated by rare, long events (L2/L3 misses), and single-thread speed must not degrade; avoid when short stalls dominate (refill cost cannot be recovered).
\item \textbf{Fine-grained} for throughput-oriented cores with many threads where single-thread latency is secondary (server/web/DB); needs per-cycle switching hardware.
\item \textbf{SMT} when you already have a wide out-of-order superscalar (renaming, large physical register set, OoO completion): \emph{the marginal hardware is small}. Not worthwhile for narrow in-order cores.
\end{itemize}
\end{dec}

\subsubsection{Why SMT is cheap on an OoO core}
Dynamically scheduled processors already have: a large set of virtual (physical) registers that can hold several threads' register sets; \textbf{register renaming} gives unique register ids so instructions from different threads mix in the datapath without confusing sources and destinations; \textbf{out-of-order completion} lets threads proceed independently. You only add a per-thread renaming table and separate PCs; independent commit is obtained by logically keeping a separate ROB per thread.

\subsubsection{SMT design challenges and costs (IBM Power5 case study)}
\begin{itemize}
\item \textbf{Single-thread performance:} SMT needs fine-grained scheduling; a \textbf{preferred-thread} approach protects one thread but loses some throughput when that thread stalls. Power5 gives hardware-controllable \emph{relative thread priority}; with balanced threads both run slower than if each owned the machine.
\item Larger register file for multiple contexts; do not lengthen the clock cycle (more candidate instructions to issue, harder choice at commit); avoid cache/TLB conflicts between threads.
\item \textbf{Why only 2 threads in Power5?} With 4, one of the shared resources (physical registers, cache, memory bandwidth) would become a bottleneck.
\item \textbf{Changes made in Power5 to support SMT:} higher associativity in the L1 I-cache and instruction translation buffers; per-thread load/store queues; larger L2 (1.92 vs.\ 1.44\,MB) and L3; separate instruction prefetch and buffering per thread; physical registers 152 $\rightarrow$ 240; larger issue queues. \textbf{Power5 core is $\approx$24\% larger than Power4's} because of SMT.
\end{itemize}

\subsubsection{SMT performance data to quote}
\begin{center}
\small
\begin{tabular}{L{6.2cm}L{9.2cm}}
\toprule
\tblhead{\textbf{Measurement} & \textbf{Result (slides)}}\\
\midrule
Pentium 4 Extreme SMT, SPECrate (same benchmark $\times$ copies) & 1.01 (SPECint\_rate), 1.07 (SPECfp\_rate) \\
Pentium 4, 26 SPEC benchmarks paired with each other (262 runs) & Speedups 0.90--1.58, \textbf{average 1.20} \\
Power5, 8-processor server & 1.23 (SPECint\_rate), 1.16 (SPECfp\_rate) with SMT \\
Power5, two copies of each application & 0.89--1.41; most gained; \textbf{FP applications had the most cache conflicts and the least gains} \\
General statement (multicore deck) & ``Enables better threading (e.g.\ up to 30\%)'' \\
\bottomrule
\end{tabular}
\end{center}
\textbf{Reading:} SMT gives a modest, workload-dependent gain (roughly 1.0--1.4$\times$) and can \emph{hurt} (speedup $<$1) when threads fight over caches. It is \emph{not} a true parallel processor: there is a single copy of each resource, and threads cannot use the same functional unit simultaneously.

\begin{ex}[Area-efficiency argument for SMT vs.\ a second core]\bs
Power5 SMT cost $\approx$24\% extra core area for $\approx$16--23\% throughput. A second core costs about 100\% extra core area but can give up to $\sim$2$\times$ throughput for fully parallel software. So SMT is roughly \emph{area-neutral} on throughput but needs little extra design and shines when threads are stall-bound; a second core wins when software is truly parallel and area/power allow.
\end{ex}
\src{ca06 ``New Approach: Multithreaded Execution'' through ``Initial Performance of SMT''; ca08 ``Simultaneous multithreading (SMT)''}

\subsection{Multicore (chip multiprocessor, CMP)}
\begin{itemize}
\item Several cores on one processor socket; they run in parallel; \textbf{within each core threads are time-sliced} as on a uniprocessor. The OS sees each core as a separate processor and schedules threads/processes across them.
\item \textbf{Why multicore:} single-core clock frequencies cannot rise further; limited ILP; many applications are already multithreaded; general trend toward parallelism.
\item A multicore is a \textbf{shared-memory MIMD multiprocessor on one chip}: different cores run different instruction streams on different data, all sharing the same memory.
\item \textbf{Applications that benefit:} database servers, web servers, compilers, multimedia, scientific, CAD/CAM; ``anything that can be threaded''. Some applications are difficult to parallelise.
\end{itemize}

\subsubsection{Multicore vs.\ SMT, and combining them}
\begin{center}
\small
\begin{tabular}{L{3.2cm}L{6.0cm}L{6.2cm}}
\toprule
\tblhead{\textbf{} & \textbf{Multicore} & \textbf{SMT}}\\
\midrule
Resources & Each core has its own copy of resources & Single copy of each resource shared by threads \\
Core size/design & Several smaller, less powerful cores; easier to design and manufacture & One large, fast superscalar core \\
Strength & Great with thread-level parallelism; independent execution & Great single-thread performance; still mostly ILP-based \\
Weakness & Does nothing for one thread; coherence traffic & Not true parallelism; shared-resource contention; $\le$30\% gain \\
Threads & One (or more) per core & 2, 4 or sometimes 8 per core (Intel ``hyper-threads'') \\
\bottomrule
\end{tabular}
\end{center}
Four combinations: single-core non-SMT (standard uniprocessor), single-core SMT, multi-core non-SMT, \textbf{multi-core with SMT} (the best-throughput option; e.g.\ SMT dual-core runs four threads concurrently).

\subsubsection{Memory hierarchy in threaded designs}
\begin{itemize}
\item SMT only $\Rightarrow$ all caches shared. Multicore $\Rightarrow$ \textbf{L1 caches private}; L2 \textbf{private in some designs, shared in others}; \textbf{main memory always shared}. Example: dual-core hyper-threaded Xeon with private L1 and \textbf{shared} L2.
\item \textbf{Shared L2:} threads share capacity dynamically and share data without duplication\bs; risk of interference. \textbf{Private L2:} isolation and simple timing\bs; duplicated data, wasted capacity if load is unbalanced\bs. In either case sharing introduces \textbf{coherence misses} (true/false sharing; Section on miss types).
\end{itemize}

\begin{pick}[Multicore + SMT]
Server/throughput workloads and mixed multitasking desktops: plenty of independent threads, area available, TLP-friendly software. Add SMT to each core if the core is a wide OoO and memory stalls are common.
\end{pick}
\begin{avoid}
Single-threaded or hard-to-parallelise applications (give one fast core instead); extremely tight power budgets that cannot afford duplicate structures or shared-cache contention.
\end{avoid}

\begin{dec}[ILP vs.\ TLP balance (slide conclusion)]
There is no universal answer: \emph{server} market $\Rightarrow$ more TLP (many cores/threads); \emph{desktop} where single-thread speed is still a primary requirement $\Rightarrow$ keep a strong ILP core but add a few cores. ``Balance of ILP and TLP decided in marketplace.''
\end{dec}

\begin{trap}[TLP traps]
\begin{itemize}
\item SMT threads cannot use the \emph{same} functional unit at the same time; it is not a doubling of resources.
\item Coarse-grained multithreading is useless when stalls are short (pipeline refill cost dominates); fine-grained slows individual threads.
\item More threads $\ne$ more speed: Power5 stops at 2 threads because shared resources saturate; SMT can give speedup $<1$ on cache-conflict-heavy (FP) code.
\item Multicore does not accelerate a single thread; the OS must have threads to schedule.
\item False sharing causes coherence misses even though data is not shared.
\end{itemize}
\end{trap}
\src{ca08 (multicore deck: ``Multi-core architectures'' through ``Designs with private L2 caches''); ca06 ``Summary'', ``Conclusions''; Cache-Misses (coherence misses)}

